\documentclass[11pt,a4paper]{article}
\usepackage[utf8]{inputenc}
\usepackage[T1]{fontenc}
\usepackage{mathptmx}
\usepackage{amsmath,amssymb}
\usepackage[margin=2.5cm]{geometry}
\usepackage{setspace}
\usepackage{hyperref}
\usepackage{graphicx}
\usepackage{lineno}
\usepackage{xcolor}
\usepackage{booktabs}
\usepackage{subcaption}

\hypersetup{
    colorlinks=true,
    linkcolor=blue,
    urlcolor=blue,
    citecolor=blue
}

\doublespacing
\linenumbers

\setlength{\parindent}{0pt}
\setlength{\parskip}{10pt}

\begin{document}

\begin{center}
{\LARGE \textbf{A non-monotonic environmental trend in Baryonic Tully--Fisher residuals: empirical evidence and a two-field phenomenological interpretation}}\\[1cm]

{\large Jonathan Edouard Slama$^{1,*}$}\\[0.5cm]

$^{1}$Metafund Research Division, Strasbourg, France\\[0.3cm]
$^{*}$Correspondence: jonathan.slama@outlook.fr\\[0.3cm]
ORCID: 0009-0002-1292-4350\\[1cm]

\textbf{Manuscript type:} Article\\
\textbf{Manuscript \#:} 2025-12-33523\\[0.5cm]

\end{center}

\newpage

%==============================================================================
% SUMMARY
%==============================================================================

\section*{Summary}

Galaxy rotation curves remain a central challenge in astrophysics, commonly interpreted through the presence of dark matter halos whose particle nature remains undetected despite extensive searches for canonical candidates including WIMPs, axions, and primordial black holes$^{1,27}$. Modified gravity theories provide partial success at galactic scales but typically predict monotonic dependencies on environmental density$^{2,3,26}$. Here we report evidence for a statistically significant non-monotonic dependence of Baryonic Tully--Fisher Relation (BTFR) residuals on environment in the SPARC sample$^{4}$. Residuals are defined throughout as the vertical offset of each galaxy from the single global BTFR fitted to the full sample, $\Delta = \log V_{\mathrm{rot}} - [a + b \log M_{\mathrm{bar}}]$; every galaxy passing the quality cuts enters the analysis, with no sigma-clipping, no outlier rejection and no three-sigma or other deviation threshold applied at any stage. The same trend is detected with complementary support from larger linewidth-based catalogues$^{5,6}$. Residuals are enhanced in both low-density voids and high-density clusters, with a minimum in intermediate environments, a pattern not captured by the simple monotonic single-field model tested here. We propose that this pattern can be phenomenologically described by a framework involving two antagonistically coupled scalar fields, preferentially interacting with gas-rich and stellar-dominated components respectively. A distinctive prediction of this framework---a sign inversion between gas-rich and gas-poor galaxy populations---is recovered at high significance in IllustrisTNG simulations and remains a target for direct observational confirmation. We stress that the present analysis tests the existence of a non-monotonic environmental trend and its qualitative consistency with a two-field description; it does not constitute a quantitative test of the framework itself, since no explicit forward model linking the field-theoretic parameters to the fitted observables is derived here. These results highlight a previously underexplored environmental trend in galaxy dynamics and provide a new observational avenue to test multi-field extensions of gravity.

%==============================================================================
% CONCEPTUAL MOTIVATION
%==============================================================================

\section*{Conceptual motivation}

This work emerged from a simple intuition: physical quantities often reveal hidden structure when expressed as ratios. We formalized this idea into a self-referential ratio between what a gravitational system expresses dynamically and what it retains structurally:
\begin{equation}
Q = \frac{\text{expressed invariant}}{\text{retained invariant}} = \frac{a_{\mathrm{eff}}}{a_{\mathrm{bar}}}
\end{equation}
where $a_{\mathrm{eff}}$ is the observed centripetal acceleration (derived from rotation velocity and radial position) and $a_{\mathrm{bar}}$ is the centripetal acceleration from the Newtonian prediction based on the baryonic mass distribution. We emphasize that this ratio is presented here only as conceptual motivation for the framework; the actual residual analysis throughout this work uses the standard BTFR formulation $\Delta = \log V_{\mathrm{rot}} - [a + b \log M_{\mathrm{bar}}]$, not a local $V^2/R$ ratio. We write the generic rotation-velocity variable as $V_{\mathrm{rot}}$ rather than $V_{\mathrm{flat}}$, because the quantity actually used is dataset-dependent: $V_{\mathrm{rot}} = V_{\mathrm{flat}}$, the asymptotic flat rotation velocity from resolved kinematic modelling, for SPARC, and an inclination-corrected W50 linewidth proxy for ALFALFA. These two velocity definitions are not equivalent$^{14,15,16}$, and we retain the distinction wherever it matters for the interpretation. See Methods for details.

We constructed a Lagrangian based on this principle and tested it on the SPARC galaxy rotation curve database$^{4}$. The prediction was simple: if a scalar field $Q$ coupled to matter density, residuals from the Baryonic Tully-Fisher Relation (BTFR) should depend \textit{monotonically} on environmental density---either increasing or decreasing, but not both.

\subsection*{The failed prediction}

The data contradicted this expectation. Rather than a monotonic trend, SPARC residuals showed elevated values at \textit{both} density extremes: galaxies in voids and galaxies in clusters rotated faster than predicted, while field galaxies at intermediate densities followed the standard relation closely.

This was not what we expected. The single-field model predicted that if $Q$ enhanced gravity in dense regions, residuals should increase monotonically with density. If $Q$ suppressed gravity, they should decrease. Within the monotonic coupling form tested here, no choice of parameters could produce elevated residuals at both extremes simultaneously.

The simple monotonic single-field model was therefore disfavoured by the data.

\subsection*{Forensic analysis: what went wrong?}

Rather than abandoning the project, we conducted a forensic analysis. What physical mechanism could produce a U-shaped pattern? This motivated a reformulation of the framework.

A single field cannot produce non-monotonic behavior with a monotonic coupling. But \textit{two} fields with \textit{opposite} environmental responses could. If one field dominates in low-density environments and another dominates in high-density environments, their combined effect could generate precisely the observed pattern.

\subsection*{The two-field reformulation}

We introduced a complementary field $R$ (for ``reliquat''---what remains after $Q$ acts) with the following properties:
\begin{itemize}
\item $Q$ couples preferentially to gas-rich components, dominating in diffuse environments
\item $R$ couples preferentially to stellar-dominated components, dominating in dense environments
\item A cross-coupling term $\lambda_{\mathrm{QR}} Q^2 R^2$, where $\lambda_{\mathrm{QR}}$ is the dimensionless coupling constant quantifying the strength of the antagonistic interaction, generates the competition between the two fields
\end{itemize}

In voids, where gas dominates, the $Q$-field enhances effective gravity, producing positive BTFR residuals. In clusters, where stars dominate, the $R$-field produces positive residuals through its opposite coupling. At intermediate densities, the two effects approximately cancel, yielding the observed minimum.

This framework made no assumptions about the fundamental origin of these fields. It was purely phenomenological---a minimal phenomenological structure capable of describing the data.

If confirmed, this environmental signature could potentially have implications beyond galaxy dynamics: it could inform dark matter search strategies, constrain cosmological models, and provide a rare low-energy window into physics typically accessible only at particle collider scales.

%==============================================================================
% DETECTION AND COMPLEMENTARY SUPPORT
%==============================================================================

\section*{Detection of the U-shaped pattern}

We analyzed 175 galaxies from the SPARC database with high-quality rotation curves and Spitzer 3.6$\mu$m photometry. For each galaxy, we computed the BTFR residual and assigned an environmental density proxy based on local galaxy counts within 2 Mpc. Note on sample sizes: throughout this work, the SPARC sample size varies slightly across analyses depending on data availability. The full BTFR fit and model-comparison analyses use $N = 181$ galaxies (all galaxies passing quality cuts $Q \leq 2$ and inclination cuts); the primary environmental-classification analysis uses $N = 175$ galaxies (those with both BTFR residuals and a documented environmental class); the 2MRS cross-matched sensitivity analysis (Supplementary Section S1) uses $N = 169$ galaxies (those additionally matched to the 2MRS catalogue with complete sky coordinates and distance estimates).

Figure 1 shows the result. A clear U-shaped pattern is detected: residuals are elevated by $+0.05$ to $+0.08$ dex in both voids and clusters, with a minimum near zero in field environments.

\begin{figure}[htbp]
\centering
\includegraphics[width=0.95\textwidth]{figure1_ushape.png}
\caption{\textbf{Detection and complementary support for the U-shaped environmental pattern.} \textbf{A}, SPARC galaxies with documented environmental classification ($N=175$) showing BTFR residuals versus environmental density. The quadratic fit yields $a = +1.33 \pm 0.25$ ($p < 0.001$). \textbf{B}, Qualitative consistency with ALFALFA ($N=21{,}834$), where $a = +0.070 \pm 0.028$ ($p = 0.0065$). \textbf{C}, Little THINGS dwarf galaxies ($N=40$) showing consistent but statistically limited signal. \textbf{D}, Cross-dataset comparison showing positive curvature estimates across all samples, with Little THINGS statistically inconclusive.}
\label{fig:ushape}
\end{figure}

We quantified the pattern by fitting a quadratic model $\Delta = a(\chi - c)^2 + d$, where $\Delta$ is the BTFR residual, $\chi$ is the environmental density proxy normalized to $[0, 1]$, $a$ is the curvature coefficient (positive for U-shape), $c$ is the location of the residual minimum, and $d$ is the baseline offset at the minimum. The best-fit curvature is $a = +1.33 \pm 0.25$ ($p < 10^{-6}$, two-tailed Student-t test on the fit coefficient against the null hypothesis $a = 0$; confirmed by permutation test with $10{,}000$ random shuffles).

To assess whether the quadratic form is required by the data rather than imposed by our fitting choice, we performed a formal model comparison on the SPARC sample using five competing functional forms: flat (constant), linear, quadratic, broken-line, and cubic spline. All non-monotonic models (quadratic, broken-line, spline) are strongly preferred over monotonic alternatives ($\Delta\mathrm{AIC} > 39$, $\Delta\mathrm{BIC} > 32$ versus linear; $\Delta\mathrm{AIC} > 53$, $\Delta\mathrm{BIC} > 42$ versus flat), indicating that the data strongly favour non-monotonic over monotonic environmental dependence within the model classes tested. Cubic spline and broken-line provide slightly better fits than pure quadratic ($\Delta\mathrm{AIC} \approx 11$, $\Delta\mathrm{BIC} \approx 7$), suggesting the underlying structure may be more complex than a pure parabola. We retain the quadratic parameterization in the main text for its physical interpretability (the curvature coefficient $a$ has a direct connection to the antagonistic Q-R coupling) while acknowledging that more flexible descriptive models can capture additional features of the residual pattern. Detailed comparison tables and visualisations are provided in Supplementary Section S2.

\subsection*{Complementary support from independent surveys}

Given the novelty of this result, we tested whether a similar qualitative trend appears in an independent survey. We extended the analysis to the ALFALFA survey$^{5}$, which provides 21,834 HI-selected galaxies with homogeneous velocity measurements across a much larger volume.

The U-shape signal is recovered in ALFALFA with $a = +0.070 \pm 0.028$ ($p = 0.0065$). We emphasize that ALFALFA W50 linewidths are integrated proxies that differ systematically from resolved rotation velocities$^{14,15,16}$, and the much smaller amplitude likely reflects this measurement systematics rather than supporting a direct quantitative comparison with SPARC. The ALFALFA analysis should therefore be interpreted as a complementary check that the non-monotonic environmental signature appears with a different velocity proxy and a much larger sample, not as a strict replication of the SPARC amplitude. Given the fundamentally different velocity definitions adopted by SPARC and ALFALFA, we phrase the comparison cautiously: the two datasets suggest a similar qualitative environmental trend (positive curvature, non-monotonic dependence on environment) rather than a matched measurement of the same quantity.

For completeness we also examined Little THINGS (40 dwarf irregulars), finding $a = +0.29 \pm 0.32$ ($p = 0.39$). The small sample precludes a statistically meaningful test, but the sign of the coefficient is consistent with the pattern observed in SPARC and ALFALFA. We do not claim this as an independent confirmation.

The U-shape is detected with high significance in SPARC ($p < 10^{-6}$, $N=175$, environmental classification based on documented galaxy-group memberships) with complementary support from ALFALFA ($p = 0.0065$, $N=21{,}834$, density-based environmental classification from 2MRS counts). The two primary datasets thus rely on \textit{complementary} environmental definitions: SPARC uses catalog-based structure membership (Virgo cluster, Ursa Major group, etc.), while ALFALFA uses local 3D galaxy density at large scales. The detection of the same qualitative U-shape in both datasets, despite their different proxies, samples, and velocity measurements, makes a purely methodological artifact less likely, although residual systematics cannot be excluded. Little THINGS shows a consistent sign but is statistically inconclusive given its size.

We emphasise that the three datasets analysed here do not share a common environmental variable. SPARC uses a catalogue-based structural classification (void / field / group / cluster), ALFALFA uses a local three-dimensional galaxy-density estimate built from 2MRS counts$^{19}$, and IllustrisTNG uses the simulation's own local density, supplemented in the Supplementary Material by the distance to the fifth nearest neighbour. These quantities are related but not equivalent, they are sensitive to different physical scales, and they need not produce identical rank orderings of galaxy environments$^{17,18}$. We have not demonstrated that they trace the same underlying physical quantity, and we do not attempt to place their fitted curvature amplitudes on a common scale. The cross-dataset comparison presented in this work should therefore be read as \textit{qualitative}: what recurs across the three analyses is the sign of the curvature and the non-monotonic shape of the environmental dependence, not a quantitatively consistent measurement of a single environmental effect. Quantitative comparison would require applying a single, homogeneously defined environmental estimator to all three samples, which is beyond the scope of the present work.

To further assess robustness, we performed a sensitivity analysis testing six alternative environmental definitions on SPARC galaxies, using galaxy counts in spheres of radii 1, 2, 3, 5, 7, and 10 Mpc computed independently from the 2MASS Redshift Survey (2MRS). This analysis reveals a physically meaningful distinction: the U-shape is strongest when environment is defined by \textbf{catalog-based structure membership} (SPARC original classification: $a = +1.33$$\pm$0.25, $\sigma = 5.26$), whereas it weakens when using \textbf{instantaneous 3D density} in fixed spheres. We interpret this as a physically meaningful distinction: the QO+R framework specifically predicts that environmental effects emerge from the structure-dependent dominance of antagonistic Q (gas-philic) and R (stellar-philic) fields, not from local density per se. A galaxy in the outskirts of a virialized cluster shares the R-dominated regime of its host structure, even if its instantaneous 3D density is comparable to a field galaxy. Detailed sensitivity tables and supporting figures are provided in Supplementary Section S1.

\textbf{Limitations of the adopted environmental proxy.} We must be explicit about what this sensitivity analysis does and does not establish, because it bears directly on how the SPARC result should be read. The primary SPARC analysis rests on the catalogue-based structure-membership classification, and the supplementary tests show that the statistical significance of the U-shape depends appreciably on which environmental definition is adopted: the signal is strongest for structure membership and weakens when environment is instead measured as an instantaneous galaxy count in fixed-radius spheres. Our interpretation---that structure membership captures the physically relevant quantity while fixed-radius density does not---is a plausible reading, but it is an interpretation, not a validation. It does not substitute for a test against the environmental estimators most widely used in the literature, namely $n$th-nearest-neighbour density, host halo mass, central/satellite classification, or tidal-interaction indicators$^{17,18,21}$, none of which is available for the full SPARC sample in a homogeneous form. Different environmental measures are known to probe genuinely different physical scales and to correlate only imperfectly with one another$^{17,18}$, so a result that is strong under one definition and weaker under another cannot yet be regarded as established independently of that choice. We therefore state plainly that the amplitude and significance of the SPARC U-shape reported here are conditional on the catalogue-based environmental classification, that its robustness across standard environmental estimators remains to be demonstrated, and that a homogeneous re-analysis using nearest-neighbour densities and halo-based classifications for a resolved-kinematics sample is a necessary next step before the trend can be considered environmentally robust.

%==============================================================================
% THE DISCRIMINATING PREDICTION
%==============================================================================

\section*{The discriminating prediction}

The two-field framework yields a prediction that is not naturally produced by the simple single-field model tested in the previous section: the U-shape should \textit{invert} between gas-dominated and stellar-dominated populations.

The logic was straightforward. If $Q$ couples to gas and produces the U-shape in gas-rich systems, then $R$-dominated systems (gas-poor, high stellar mass) should show the \textit{opposite} pattern---an inverted U-shape with $a < 0$.

This was not a post-hoc adjustment. It was a prediction derived from the framework's structure before we examined the relevant data. Direct observational tests of this sign inversion in gas-poor early-type galaxies are currently challenging due to the requirement of high-quality kinematic data in low-density environments, but preliminary indications from MaNGA and SAMI integral-field-unit surveys are consistent with the prediction and will be the subject of a forthcoming companion paper. The present analysis uses IllustrisTNG as a controlled environment to test whether the sign inversion can be reproduced by standard $\Lambda$CDM physics alone.

\subsection*{Testing the prediction}

We used the IllustrisTNG simulations$^{9,20}$, which provide gas and stellar masses for over 600,000 galaxies with known environments. To test this prediction rigorously, we adopted a stratified design: galaxies are first grouped by \textbf{internal properties} (gas fraction and stellar mass) to identify populations where the QO+R framework predicts Q-dominance (gas-rich) versus R-dominance (gas-poor, stellar-massive). Within each population, the U-shape coefficient is then fitted against an \textbf{environmental variable}: the local 3D galaxy density (standardised, computed from the simulation's spatial information within a fixed-radius sphere). The internal properties thus serve as \textbf{covariates} that select the predicted Q/R regime, while the dependence on environment is tested using a direct 3D density proxy. A supplementary analysis (Section S3) shows that the sign inversion is strongest in dense environments and disappears in isolated systems when using a direct 3D isolation proxy (distance to the 5th nearest neighbour).

\begin{itemize}
\item Gas-rich ($f_{\mathrm{gas}} > 0.5$, $N = 444,374$): $a = +0.017 \pm 0.008$
\item Gas-poor + high $M_*$ ($f_{\mathrm{gas}} < 0.2$, $\log M_* > 10.5$, $N = 8{,}779$): $a = -0.014 \pm 0.001$. Note: the gas-poor high-mass sample used in this main TNG density-binned analysis contains $N = 8{,}779$ galaxies after binning and quality cuts, while the isolation-controlled supplementary analysis (Section S3) uses $N = 9{,}386$ R-dominated galaxies before the final density-binning cuts.
\end{itemize}

\begin{figure}[htbp]
\centering
\includegraphics[width=0.95\textwidth]{figure2_inversion.png}
\caption{\textbf{Sign inversion test of the two-field prediction in TNG300 simulations.} \textbf{A}, U-shape coefficient $a$ as a function of gas fraction in TNG300 ($N=623{,}609$): Q-dominated (gas-rich) populations show $a > 0$ (standard U-shape), while R-dominated (gas-poor) populations show $a < 0$ (inverted U-shape). \textbf{B}, Gas-poor galaxies split by stellar mass: Low-mass ($N=115{,}980$, $a=+0.052\pm0.012$), mid-mass ($N=25{,}665$, $a\approx 0$), high-mass ($N=8{,}779$, $a=-0.014\pm0.003$). \textbf{C}, Detailed breakdown across 5 R-dominance categories: standard U-shape recovered in Q-dominated/gas-rich/low-mass; inverted U-shape ($a<0$) in mid-mass, high-mass, and extreme R-dominated quiescent systems. The sign difference is highly significant within the TNG300-1 simulated sample. The isolation-controlled R-dominated analysis is provided in Supplementary Section S3; broader category definitions are documented in the companion Zenodo deposit (Slama 2025). This is a simulation-based test of the discriminating prediction, not an observational confirmation; observational follow-up in gas-poor early-type galaxies will be the subject of a forthcoming companion paper.}
\label{fig:inversion}
\end{figure}

The sign difference is highly significant within the TNG300-1 simulated sample (Figure 2). Because the galaxies are drawn from a $\Lambda$CDM simulation with baryonic subgrid prescriptions, this significance quantifies the internal detection of the pattern in the simulated catalogue, not an observational confirmation of the QO+R framework. Gas-rich populations show the standard U-shape; gas-poor, high-mass populations show the inverted pattern. Detailed isolation-controlled statistics for the R-dominated population are provided in Supplementary Section S3.

\textbf{Environment as the primary variable: a multivariate test.} The stratified design above has an intrinsic limitation, which we now address directly. Because the Q- and R-dominated populations are defined using \textit{internal} galaxy properties (gas fraction and stellar mass) \textit{before} their environmental dependence is examined, environmental effects and population-dependent differences associated with morphology, quenching history or baryonic subgrid physics are not cleanly separated by that design alone. To disentangle them we repeated the analysis with environment as the primary independent variable and the internal properties demoted to statistical controls, using the TNG100-1 catalogue at $z=0$ ($N = 53{,}363$ galaxies with $M_* > 10^9 M_\odot$). We fitted a hierarchy of nested ordinary-least-squares models to the individual BTFR residuals, with all predictors standardised and heteroskedasticity-robust (HC3) standard errors: (M0) residual as a quadratic function of local environmental density alone; (M1) the same, adding gas fraction and $\log M_*$ as linear controls; and (M2) the same, adding the interaction terms $\rho \times f_{\mathrm{gas}}$, $\rho \times \log M_*$ and $\rho^2 \times f_{\mathrm{gas}}$.

The quadratic environmental coefficient remains positive and statistically significant at every level of the hierarchy: $+0.079 \pm 0.002$ in M0, $+0.029 \pm 0.002$ in M1, and $+0.0095 \pm 0.0017$ in M2 (all $p < 10^{-7}$). We draw attention to the amplitude rather than the significance: controlling for gas fraction and stellar mass attenuates the curvature by roughly a factor of three, and adding the interaction terms attenuates it by a further factor of three, so that the coefficient in the fully controlled model is about eight times smaller than in the environment-only fit. A substantial part of the raw environmental curvature is therefore attributable to the covariance between environment and internal galaxy properties, and only a residual---though robustly non-zero---component is associated with environment once those properties are held fixed.

The interaction structure of M2 also allows the sign inversion to be examined without any prior subsample selection. The marginal environmental slope $\partial \Delta / \partial \rho$ evaluated at fixed stellar mass changes sign as a continuous function of gas fraction, from $-0.077 \pm 0.002$ for gas-rich systems ($f_{\mathrm{gas}}$ one standard deviation above the mean) to $+0.141 \pm 0.006$ for gas-poor systems ($f_{\mathrm{gas}}$ one standard deviation below the mean), with an interaction coefficient $\rho \times f_{\mathrm{gas}} = -0.109 \pm 0.003$. The inversion of the environmental response with gas content is thus recovered in a purely multivariate treatment in which no population is pre-selected. Full regression tables and a three-panel summary figure are provided in Supplementary Section S4, and the analysis script is included in the code deposit.

We nonetheless retain the caveat in explicit form. This multivariate treatment controls for gas fraction and stellar mass, but it does not control for every galaxy property that co-varies with environment in a hydrodynamical simulation---morphology, assembly and quenching history, black-hole feedback state and satellite status among them---nor does it establish causality. The attenuation of the curvature under controls should be read as evidence that a meaningful fraction of the raw signal is population-driven, and the surviving component as a lower bound on any genuinely environmental effect within TNG.

This behaviour is difficult to accommodate within the simple monotonic single-field model considered here, and therefore provides a discriminating target for the two-field phenomenological framework. However, it should not be interpreted as excluding all possible single-field, baryonic, or population-mixing alternatives.

%==============================================================================
% ROBUSTNESS
%==============================================================================

\section*{Statistical robustness}

Given the novelty of this pattern, we subjected the results to a series of statistical resampling tests (Figure 3). We state at the outset what these tests do and do not cover, since the distinction matters for their interpretation: they assess the \textit{internal statistical robustness} of the fitted trend---its stability under resampling, perturbation and permutation of the existing measurements---and nothing more. They do not address observational or astrophysical systematic uncertainties, including environmental misclassification, distance uncertainties, inclination errors and uncertainties in baryonic mass estimates, none of which is propagated in the analysis below. Describing these tests as an exhaustive assessment of robustness would overstate what has been demonstrated.

\begin{figure}[htbp]
\centering
\includegraphics[width=0.95\textwidth]{figure3_robustness.png}
\caption{\textbf{Statistical robustness of the U-shape detection in SPARC.} \textbf{A}, Monte Carlo perturbation (1000 iterations with realistic measurement errors): 100\% of resampled datasets preserve positive curvature; mean coefficient $a = 1.33$. \textbf{B}, Bootstrap resampling (1000 samples): 95\% confidence interval $[0.94, 1.75]$, excluding zero. \textbf{C}, Jackknife leave-one-environment-out: signal persists when removing any single environmental bin (3/4 environments yield $a > 0$ after removal). \textbf{D}, Cross-validation comparing linear vs quadratic: quadratic preferred with 7.9\% RMSE improvement. \textbf{E}, Permutation testing: observed coefficient exceeds 99.99\% of the null distribution ($p < 0.0001$). These five tests assess internal statistical robustness only: they disfavour explanations based solely on outliers or random statistical fluctuation, but they do not propagate observational systematics (environmental misclassification, distance and inclination uncertainties, baryonic mass errors), which remain unquantified.}
\label{fig:robustness}
\end{figure}

\begin{itemize}
\item \textbf{Monte Carlo perturbation} (1000 iterations): 100\% of samples preserve positive curvature
\item \textbf{Bootstrap resampling}: 95\% CI = [0.94, 1.75], excluding zero
\item \textbf{Jackknife}: Signal persists when removing any single environmental bin
\item \textbf{Cross-validation}: Quadratic model preferred over linear (7.9\% RMSE improvement)
\item \textbf{Permutation testing}: $p < 0.0001$ against random shuffling
\end{itemize}

These tests collectively disfavour explanations of the pattern in terms of outliers or statistical fluctuation alone. They do not, however, constrain observational and methodological systematics, which remain the dominant unquantified uncertainty in this analysis. In particular, misclassification of environments, systematic distance errors, inclination-dependent velocity biases and systematic uncertainties in stellar mass-to-light ratios could each in principle induce or modulate a curvature in the residual--environment relation, and none is excluded by the resampling tests reported here. A full systematic error budget would require a homogeneous re-analysis with propagated distance, inclination and mass-modelling uncertainties, which we identify as necessary future work rather than something achieved in the present manuscript.

%==============================================================================
% SCREENING AND PRECISION TESTS
%==============================================================================

\section*{Screening and precision tests}

Any modification to gravity must satisfy stringent constraints from Solar System observations$^{10}$ and laboratory experiments$^{11}$. Several screening mechanisms have been proposed to reconcile these constraints with cosmologically relevant modifications, notably the Chameleon$^{3}$, Symmetron$^{22}$ and Vainshtein$^{23}$ mechanisms. How does the framework evade these bounds while producing detectable effects at galactic scales?

The answer lies in environmental screening. The screening mechanism is modeled via a density-dependent suppression function:
\begin{equation}
\mu(\rho) = \exp(-\rho/\rho_0)
\end{equation}
where $\rho_0$ is the characteristic density scale. This suppresses modifications to gravity in high-density regions ($\rho \gg \rho_0$), by construction ensuring compatibility with local tests; the exponential form and the value of $\rho_0$ are both adopted rather than derived (see below). At densities characteristic of the Solar System or stellar interiors, the scalar field effects are naturally suppressed. The framework predicts null results in high-density environments---and the data are consistent with this expectation in the present analysis.

The characteristic density $\rho_0$ defines the transition scale between the screened (high-density) and unscreened (low-density) regimes. We must be unambiguous about the status of this quantity. The value $\rho_0 \sim 10^{-24}$ g cm$^{-3}$ is \textit{not} predicted by the framework. It is an adopted phenomenological parameter, fixed \textit{a posteriori} by requiring that modifications remain active at galactic densities while being suppressed at Solar System densities---that is, it is chosen to reproduce the required screening behaviour rather than derived from it. The framework as formulated here contains no independent determination of $\rho_0$ from its own field content, and the agreement between the adopted value and the observed galactic/Solar-System separation is therefore a consistency requirement that has been imposed, not a successful prediction that has been tested. This is qualitatively similar to the Chameleon mechanism$^{3}$, with the difference that our framework involves two coupled fields rather than a single one. A rigorous derivation of $\rho_0$ from an underlying field-theoretic action, including the explicit equations of motion for the coupled Q and R fields in the presence of a matter source, is the subject of a companion theoretical paper currently in preparation. We acknowledge that the present treatment is phenomenological in nature: the exponential form $\mu(\rho) = \exp(-\rho/\rho_0)$ is motivated by analogy with the broader class of density-dependent screening mechanisms (Chameleon, Symmetron, Vainshtein) but is not derived from first principles in this manuscript.

\begin{figure}[htbp]
\centering
\includegraphics[width=0.95\textwidth]{figure4_screening.png}
\caption{\textbf{Consistency with precision tests through environmental screening.} \textbf{A}, Scale hierarchy showing the QO+R active region ($\lambda_{\mathrm{QR}} \sim 1$ kpc) versus precision test scales (Eot-Wash, LLR, Cassini-Saturn, galactic). \textbf{B}, Observational constraints vs QO+R predictions for Eot-Wash, PPN-Cassini, Lunar Laser Ranging, and Mercury precession; all predictions lie well below current observational limits. \textbf{C}, Screening mechanism: effective coupling $\alpha_{\mathrm{eff}}$ as a function of local density, showing strong suppression for $\rho \gtrsim 10^{-24}$ g cm$^{-3}$ (Solar System and stellar regimes). The galactic regime ($\rho \sim 10^{-25}$ to $10^{-26}$ g cm$^{-3}$) lies in the active region. Key insight: QO+R effects emerge at galactic scales where density gradients are large and environmental variations span void-cluster; at Solar System densities, the framework recovers standard General Relativity.}
\label{fig:screening}
\end{figure}

Specific null results consistent with screening (Figure 4):
\begin{itemize}
\item Globular clusters (Baumgardt \& Hilker catalogue$^{24}$): $r = +0.12$, $p = 0.30$ (no correlation)
\item Wide binaries (Gaia DR3$^{25}$): No significant deviation from Newtonian dynamics
\item Solar System (Cassini): $|\gamma_{\mathrm{PPN}} - 1| < 2.3 \times 10^{-5}$
\end{itemize}

The framework is not in obvious conflict with the precision tests considered here, while predicting detectable effects only in diffuse astrophysical environments. A full constraint analysis across other regimes remains necessary.

%==============================================================================
% THEORETICAL CONNECTION
%==============================================================================

\section*{Theoretical connection}

We then asked whether this phenomenological structure could connect to fundamental physics. The two-field Lagrangian is structurally compatible with the moduli sector of string compactification scenarios$^{7,8}$. The $Q$ field may be identified with moduli controlling gauge couplings, while $R$ corresponds to geometric moduli. The cross-coupling $\lambda_{\mathrm{QR}}$ emerges from the K\"ahler potential structure.

Moduli stabilisation mechanisms generically produce dimensionless cross-couplings of order unity without fine-tuning, which is the sense in which the two-field structure is compatible with that setting.

\textbf{What is and is not tested here.} We wish to be entirely explicit about the epistemic status of this section, because it determines how the empirical results of this paper should be read. No explicit forward model is derived in this work that maps the parameters of the two-field framework---the coupling strength $\lambda_{\mathrm{QR}}$, the screening parameters, or the field amplitudes---onto the fitted observables, that is, onto the curvature coefficient $a$ of the residual--environment relation and its dependence on gas fraction. Without such a mapping, the fitted curvature cannot be converted into a measurement of any theoretical quantity, and conversely no value of any theoretical quantity predicts a specific value of the curvature. It follows that the observational analysis presented in this manuscript tests the \textit{presence of a non-monotonic environmental trend} and the \textit{qualitative sign structure} of its dependence on gas content; it does not test the two-field framework quantitatively, and it does not measure the coupling between the fields.

Accordingly, the present work should be understood as a \textbf{phenomenological consistency test}: the observed pattern is of the kind that a two-field antagonistic structure would produce, and is not of the kind that the simple monotonic single-field model tested here produces. That is a meaningful but limited statement. Any reading of the curvature coefficient as direct evidence for a Q--R interaction, or of its stability across samples as a measurement of a fundamental coupling, would go beyond what the analysis supports. The connection to moduli physics is offered as a possible interpretive context, not as a proof, not as a derivation, and not as a validated identification; deriving the forward model that would make this connection quantitatively testable is the necessary next step, and it has not been carried out here.

%==============================================================================
% FALSIFIABLE PREDICTIONS
%==============================================================================

\section*{Falsifiable predictions}

The framework makes specific predictions testable with upcoming surveys. Table 1 summarizes the key predictions and their observational tests.

\begin{table}[htbp]
\centering
\caption{\textbf{Summary of falsifiable predictions and observational tests.}}
\small
\begin{tabular}{@{}p{0.26\textwidth}p{0.34\textwidth}p{0.30\textwidth}@{}}
\toprule
\textbf{Prediction} & \textbf{Observational test} & \textbf{Mission/Instrument} \\
\midrule
U-shape in UDGs & Velocity dispersion / resolved kinematics & Euclid + ground-based spectroscopy, SKA pathfinders \\
\addlinespace
Sign inversion (gas-poor) & Rotation curves in voids & SKA \\
\addlinespace
Shear excess in filaments & Weak lensing & LSST, Euclid \\
\addlinespace
Population dependence & Environment vs gas fraction & 4MOST, WEAVE \\
\bottomrule
\end{tabular}
\label{tab:predictions}
\end{table}

\begin{enumerate}
\item \textbf{Population dependence}: DESI and Euclid will measure environments for millions of galaxies. The U-shape amplitude should correlate with gas fraction, with sign inversion for gas-poor populations.

\item \textbf{Ultra-diffuse galaxies}: Systems with extremely low surface brightness should show enhanced velocity dispersions by $\sim$6\% in diffuse environments.

\item \textbf{Void galaxies}: SKA pathfinders will provide rotation curves for void galaxies, which should show systematically positive BTFR residuals.

\item \textbf{Weak lensing}: Filamentary structures should show 3--5\% excess shear compared to GR predictions.
\end{enumerate}

The framework would be \textbf{falsified} if:
\begin{itemize}
\item The U-shape curvature becomes negative in any large gas-rich sample
\item The sign inversion disappears in independent surveys
\item The coupling shows strong scale dependence outside the range 0.5--2.0
\item High-density systems show non-null signals above screening predictions
\end{itemize}

%==============================================================================
% DISCUSSION
%==============================================================================

\section*{Discussion}

This work began with a failed prediction and led to the identification of a previously underexplored non-monotonic environmental trend. The simple monotonic single-field hypothesis was disfavoured by SPARC data, but forensic analysis revealed an unexpected U-shaped pattern. This pattern motivated a two-field phenomenological interpretation involving antagonistically coupled fields, which suggested a discriminating qualitative signature (sign inversion) that is recovered in IllustrisTNG simulations, including under a multivariate treatment in which environment is the primary variable and internal galaxy properties are controlled. Observational confirmation in gas-poor early-type galaxies remains a key target for future work.

The conceptual development exemplifies how falsification drives progress. Had we abandoned the project after the initial failure, the U-shape would have remained hidden. Instead, asking ``what went wrong?'' led to a richer framework with stronger predictive power.

Several caveats deserve emphasis:
\begin{itemize}
\item Environmental density estimates carry systematic uncertainties
\item The sign inversion is recovered in simulations; observational confirmation with real gas-poor samples is needed
\item Alternative astrophysical explanations cannot be definitively excluded
\end{itemize}

To these we add the three limitations identified above, which we consider the most consequential. First, the amplitude and significance of the SPARC signal are conditional on the catalogue-based environmental classification and have not been validated against standard environmental estimators. Second, the environmental variables used in SPARC, ALFALFA and IllustrisTNG are not equivalent, so the agreement between datasets is qualitative in nature. Third, and most importantly, no forward model links the parameters of the two-field framework to the fitted observables, so this work constitutes a phenomenological consistency test rather than a quantitative test of the framework; the curvature coefficient should not be read as a measurement of an underlying coupling.

With those limitations stated, the combination of a statistically robust SPARC detection, a qualitatively similar trend in ALFALFA, a simulation-based test of the sign inversion that survives multivariate control for gas fraction and stellar mass, and the absence of conflict with precision constraints suggests that the observed environmental pattern warrants further investigation---primarily through homogeneous environmental re-analysis and through the derivation of a forward model capable of turning the framework into a quantitatively falsifiable one.

These results highlight a previously underexplored environmental trend in galaxy dynamics and provide a new observational avenue to test multi-field extensions of gravity. Future surveys will determine whether the two-field interpretation survives continued scrutiny.

%==============================================================================
% METHODS
%==============================================================================

\section*{Methods}

\subsection*{SPARC analysis}
We used SPARC v2.1$^{4}$ (Lelli, McGaugh \& Schombert 2016), selecting galaxies with quality flag $Q \leq 2$ and inclination $30^\circ < i < 80^\circ$. Baryonic masses follow the standard prescription: $M_{\mathrm{bar}} = M_* + 1.33 M_{\mathrm{HI}}$. The Baryonic Tully-Fisher Relation$^{12,13,14,15,16}$ was fitted to the full SPARC sample as $\log V_{\mathrm{rot}} = a + b \log M_{\mathrm{bar}}$ (with $V_{\mathrm{rot}} = V_{\mathrm{flat}}$ for SPARC) via ordinary linear least-squares regression in log-log space, yielding $a = -1.058$ and $b = 0.344$ (corresponding to slope $\alpha = 1/b \approx 2.91$, consistent with the slope range $\alpha = 2.9$--$4.0$ reported in the literature for different velocity definitions and sample selections). \textit{Note on fit orientation:} we adopt the orientation $\log V_{\mathrm{rot}}$ as a function of $\log M_{\mathrm{bar}}$, treating $M_{\mathrm{bar}}$ as the independent variable. This choice reflects the analysis goal: we study how environmental factors modify the velocity prediction for galaxies of a given baryonic mass. Environmental classifications are taken from documented group and cluster memberships, and the limitations of this choice are discussed explicitly in the main text. The reverse orientation ($\log M_{\mathrm{bar}}$ as a function of $\log V_{\mathrm{rot}}$), more commonly used when probing fundamental physics from the BTFR slope itself, yields a numerically different slope due to bivariate scatter, but we verified that it does not change the qualitative environmental trend reported here (the sign and statistical significance of the quadratic curvature coefficient are preserved). Residuals are defined as the vertical offset of each galaxy from this single global fit, $\Delta = \log V_{\mathrm{rot}} - [a + b \log M_{\mathrm{bar}}]$, evaluated for every galaxy passing the quality and inclination cuts. No sigma-clipping, no outlier rejection and no deviation threshold of any kind (in particular no three-sigma criterion) is applied at any stage: the full distribution of residuals enters the environmental analysis, including galaxies lying far from the relation. The residual is therefore a continuous per-galaxy quantity, not a flag identifying discrepant objects. Environmental densities are derived from NED galaxy counts within a 2 Mpc projected radius around each SPARC galaxy.

\subsection*{ALFALFA analysis}
We used the ALFALFA $\alpha.100$ catalog$^{5}$, selecting sources with quality codes 1--2. Inclination-corrected W50 linewidths were used as the integrated velocity proxy $V_{\mathrm{rot}}$ for this sample (note that W50 widths systematically differ from resolved $V_{\mathrm{flat}}$$^{14,15,16}$, as discussed in the main text). Environmental densities were estimated from 2MASS Redshift Survey$^{19}$ galaxy counts on large scales, providing a density-based environmental classification complementary to the catalog-based SPARC structure-membership classification.

\subsection*{IllustrisTNG analysis}
TNG100-1 and TNG300-1$^{9,20}$ at $z = 0$. Selection: $M_* > 10^9 M_\odot$. Gas fractions: $f_{\mathrm{gas}} = M_{\mathrm{gas}}/(M_{\mathrm{gas}} + M_*)$. Local environmental density is computed from the simulation's three-dimensional spatial information within a fixed-radius sphere and used in logarithmic form; the supplementary isolation analysis uses the distance to the fifth nearest neighbour.

\textit{Multivariate analysis.} For the analysis in which environment is treated as the primary independent variable, we used the TNG100-1 catalogue at $z=0$ ($N = 53{,}363$). Each galaxy contributes its individual BTFR residual, computed exactly as for the observational samples. Predictors (log environmental density, gas fraction, $\log M_*$) were standardised to zero mean and unit variance, so that fitted coefficients are directly comparable in magnitude. Three nested ordinary-least-squares models were fitted---environment only (quadratic), plus linear controls, plus interaction terms---using heteroskedasticity-consistent (HC3) standard errors throughout, with model preference assessed by AIC and BIC. Marginal environmental slopes at specified gas fractions were obtained from the full parameter covariance matrix. The complete regression tables, the analysis script and the summary figure are provided in Supplementary Section S4 and in the code deposit.

\subsection*{Statistical methods}
Quadratic fitting with weighted least squares. Significance via t-statistics and permutation tests (10,000 iterations). Bootstrap confidence intervals (1000 samples). Multivariate regressions use ordinary least squares with HC3 heteroskedasticity-consistent standard errors. All quoted uncertainties are statistical; observational and astrophysical systematic uncertainties are not propagated in any of these tests, as discussed in the section on statistical robustness.

\subsection*{Code availability}
All custom analysis code used in this study---the BTFR fitting and residual computation, the environmental sensitivity analysis, the statistical resampling tests, and the multivariate IllustrisTNG regression---has been deposited in the Zenodo repository, which assigns a permanent DOI: \url{https://doi.org/10.5281/zenodo.17806441}. The deposit contains the analysis scripts, the derived data tables required to reproduce every figure and table in this manuscript and its Supplementary Information, and a README documenting the execution order and software environment (Python 3, NumPy, SciPy, pandas, statsmodels, Matplotlib). The development repository is additionally available at \url{https://github.com/JonathanSlama/QO-R-JEDSLAMA}; the Zenodo deposit is the citable archived version of record for this paper.

\subsection*{Data availability}
All data used in this study are publicly available from the following sources: SPARC database (Lelli et al., 2016, Astron. J. 152, 157), ALFALFA $\alpha.100$ survey (Haynes et al., 2018, Astrophys. J. 861, 49), Little THINGS (Hunter et al., 2012, Astron. J. 144, 134), IllustrisTNG simulations (available at \url{https://www.tng-project.org/}), and the 2MASS Redshift Survey (Huchra et al., 2012, Astrophys. J. Suppl. 199, 26). All analysis code and processed datasets are archived on Zenodo under a permanent DOI (see Code availability) and are additionally available on GitHub (\url{https://github.com/JonathanSlama/QO-R-JEDSLAMA}).

%==============================================================================
% REFERENCES
%==============================================================================

\section*{References}

\begin{enumerate}
\item Aprile, E. \textit{et al.} Dark matter search results from a one ton-year exposure of XENON1T. \textit{Phys. Rev. Lett.} \textbf{121}, 111302 (2018).

\item Milgrom, M. A modification of the Newtonian dynamics as a possible alternative to the hidden mass hypothesis. \textit{Astrophys. J.} \textbf{270}, 365--370 (1983).

\item Khoury, J. \& Weltman, A. Chameleon cosmology. \textit{Phys. Rev. D} \textbf{69}, 044026 (2004).

\item Lelli, F., McGaugh, S. S. \& Schombert, J. M. SPARC: Mass models for 175 disk galaxies. \textit{Astron. J.} \textbf{152}, 157 (2016).

\item Haynes, M. P. \textit{et al.} The ALFALFA extragalactic HI source catalog. \textit{Astrophys. J.} \textbf{861}, 49 (2018).

\item de Jong, J. T. A. \textit{et al.} The Kilo-Degree Survey. \textit{Astron. Astrophys.} \textbf{604}, A134 (2017).

\item Giddings, S. B., Kachru, S. \& Polchinski, J. Hierarchies from fluxes in string compactifications. \textit{Phys. Rev. D} \textbf{66}, 106006 (2002).

\item Kachru, S., Kallosh, R., Linde, A. \& Trivedi, S. P. De Sitter vacua in string theory. \textit{Phys. Rev. D} \textbf{68}, 046005 (2003).

\item Pillepich, A. \textit{et al.} First results from the IllustrisTNG simulations. \textit{Mon. Not. R. Astron. Soc.} \textbf{475}, 648--675 (2018).

\item Bertotti, B., Iess, L. \& Tortora, P. A test of general relativity using radio links with the Cassini spacecraft. \textit{Nature} \textbf{425}, 374--376 (2003).

\item Adelberger, E. G., Heckel, B. R. \& Nelson, A. E. Tests of the gravitational inverse-square law. \textit{Annu. Rev. Nucl. Part. Sci.} \textbf{53}, 77--121 (2003).

\item Tully, R. B. \& Fisher, J. R. A new method of determining distances to galaxies. \textit{Astron. Astrophys.} \textbf{54}, 661--673 (1977).

\item McGaugh, S. S., Schombert, J. M., Bothun, G. D. \& de Blok, W. J. G. The Baryonic Tully-Fisher Relation. \textit{Astrophys. J. Lett.} \textbf{533}, L99--L102 (2000).

\item Lelli, F., McGaugh, S. S., Schombert, J. M. \& Pawlowski, M. S. The baryonic Tully-Fisher relation for different velocity definitions and implications for galaxy angular momentum. \textit{Mon. Not. R. Astron. Soc.} \textbf{484}, 3267--3278 (2019).

\item Ponomareva, A. A. \textit{et al.} MIGHTEE-HI: the baryonic Tully-Fisher relation over the last billion years. \textit{Mon. Not. R. Astron. Soc.} \textbf{508}, 1195--1208 (2021).

\item Biswas, P., Kalinova, V., Roy, N., Patra, N. N. \& Tyulneva, N. The GMRT archive atomic gas survey -- II. Mass modelling and dark matter halo properties across late-type spirals. \textit{Mon. Not. R. Astron. Soc.} \textbf{524}, 6213--6228 (2023).

\item Muldrew, S. I. \textit{et al.} Measures of galaxy environment -- I. What is ``environment''? \textit{Mon. Not. R. Astron. Soc.} \textbf{419}, 2670--2682 (2012).

\item Haas, M. R., Schaye, J. \& Jeeson-Daniel, A. Disentangling galaxy environment and host halo mass. \textit{Mon. Not. R. Astron. Soc.} \textbf{419}, 2133--2146 (2012).

\item Huchra, J. P. \textit{et al.} The 2MASS Redshift Survey---description and data release. \textit{Astrophys. J. Suppl. Ser.} \textbf{199}, 26 (2012).

\item Nelson, D. \textit{et al.} The IllustrisTNG simulations: public data release. \textit{Comput. Astrophys. Cosmol.} \textbf{6}, 2 (2019).

\item Peng, Y.-J. \textit{et al.} Mass and environment as drivers of galaxy evolution in SDSS and zCOSMOS and the origin of the Schechter function. \textit{Astrophys. J.} \textbf{721}, 193--221 (2010).

\item Hinterbichler, K. \& Khoury, J. Symmetron fields: screening long-range forces through local symmetry restoration. \textit{Phys. Rev. Lett.} \textbf{104}, 231301 (2010).

\item Vainshtein, A. I. To the problem of nonvanishing gravitation mass. \textit{Phys. Lett. B} \textbf{39}, 393--394 (1972).

\item Baumgardt, H. \& Hilker, M. A catalogue of masses, structural parameters and velocity dispersion profiles of 112 Milky Way globular clusters. \textit{Mon. Not. R. Astron. Soc.} \textbf{478}, 1520--1557 (2018).

\item Gaia Collaboration. Gaia Data Release 3: summary of the content and survey properties. \textit{Astron. Astrophys.} \textbf{674}, A1 (2023).

\item Famaey, B. \& McGaugh, S. S. Modified Newtonian Dynamics (MOND): observational phenomenology and relativistic extensions. \textit{Living Rev. Relativ.} \textbf{15}, 10 (2012).

\item Bullock, J. S. \& Boylan-Kolchin, M. Small-scale challenges to the $\Lambda$CDM paradigm. \textit{Annu. Rev. Astron. Astrophys.} \textbf{55}, 343--387 (2017).
\end{enumerate}

\section*{Acknowledgements}
Data from SPARC (Lelli, McGaugh, Schombert), ALFALFA (Haynes et al.), and IllustrisTNG (Pillepich et al.).

\section*{Author contributions}
J.E.S. conceived the framework, performed all analyses, and wrote the manuscript.

\section*{Funding}
This research did not receive any specific grant from funding agencies in the public, commercial, or not-for-profit sectors.

\section*{Competing interests}
The author declares no competing interests.

\end{document}
